\chapter{GIỚI THIỆU TỔNG QUAN}
\label{chap:intro}

\section{Bối cảnh và Sự cần thiết của đề tài}
Trong bối cảnh nền công nghiệp 4.0 và hướng tới công nghiệp 5.0, xu hướng tự động hóa không chỉ dừng lại ở các dây chuyền sản xuất độc lập mà đang chuyển dịch mạnh mẽ sang sự kết hợp giữa trí tuệ thích ứng của con người và độ chính xác của máy móc, hay còn gọi là Tương tác cộng tác Người - Máy (Human-Robot Collaboration - HRC) \cite{kawshan2026digital} [1]. Mặc dù robot công nghiệp hiện đại có khả năng tự chủ ngày càng cao, sự can thiệp của con người thông qua điều khiển từ xa (teleoperation) vẫn là giải pháp thiết yếu và an toàn nhất khi xử lý các môi trường không có cấu trúc, độc hại, hoặc các tình huống khẩn cấp nằm ngoài khả năng xử lý của các thuật toán tự động \cite{torrejon2026teleoperation} [2]. 

Theo phương thức truyền thống, hệ thống điều khiển từ xa phụ thuộc nhiều vào các thiết bị đầu vào 2D như màn hình, bàn phím hoặc tay cầm (joystick) \cite{hetrick2020comparing} [3]. Những giao diện 2D này đòi hỏi người vận hành phải liên tục chuyển đổi dữ liệu hình ảnh phẳng thành các quyết định trong không gian 3D, từ đó làm tăng tải nhận thức (cognitive load), gây mệt mỏi tinh thần và dễ dẫn đến sai sót \cite{chen2023comparing} [4]. Thực tế ảo (VR - Virtual Reality) đã vươn lên như một công nghệ đột phá, cung cấp môi trường nhập vai 3D, khôi phục nhận thức chiều sâu và cho phép tương tác trực quan bằng các cử chỉ tự nhiên của con người \cite{zhang2023research} [5].

\section{Động lực nghiên cứu (Motivation)}
Động lực chính của nghiên cứu này xuất phát từ nhu cầu giải quyết "điểm nghẽn" (bottleneck) trong các hệ thống HRC hiện tại: sự thiếu đồng bộ thời gian thực và sự hạn chế của các giao diện tương tác cũ. Khi ứng dụng VR vào điều khiển robot, một trong những thách thức lớn nhất là duy trì độ trễ thấp; nếu độ trễ vượt quá ngưỡng cho phép (khoảng 150ms đến 250ms), người vận hành sẽ bị mất phương hướng, chóng mặt (cybersickness) và độ chính xác của thao tác sẽ giảm sút nghiêm trọng \cite{torrejon2026teleoperation} [6]. 

Để giải quyết vấn đề độ trễ này, công nghệ Bản sao số (Digital Twin) đã được chứng minh là một giải pháp tối ưu. Thay vì chờ đợi phản hồi vật lý từ robot ở xa, người vận hành có thể tương tác trực tiếp với Bản sao số trong môi trường VR với độ trễ gần như bằng không \cite{autolab2021telesurgery} [7]. Tuy nhiên, phần lớn các hệ thống Digital Twin hiện nay vẫn dựa trên nền tảng Robot Operating System thế hệ đầu (ROS1), vốn bộc lộ nhiều điểm yếu về quản lý băng thông và không đảm bảo tốt tính thời gian thực khi truyền tải các khối dữ liệu lớn như hình ảnh hay Point Cloud \cite{kawshan2026digital} [8]. Việc thiếu một kiến trúc phần mềm tận dụng sức mạnh của ROS2 (dựa trên DDS) kết hợp với các thiết bị VR thế hệ mới (như Meta Quest 3) chính là động lực thôi thúc sự ra đời của luận văn này.

\section{Các thách thức hiện tại (Challenges)}
Mặc dù tiềm năng là rất lớn, việc xây dựng một hệ thống điều khiển robot qua VR vẫn phải đối mặt với các thách thức học thuật và kỹ thuật sau:
\begin{itemize}
    \item \textbf{Khoảng cách Sim-to-Real và Độ trễ mạng:} Việc đồng bộ hóa giữa hai môi trường vật lý (ROS/Controller) và mô phỏng (Unity) gặp khó khăn do độ trễ truyền dẫn và xử lý. Phân tích cho thấy phần cứng thường là nút thắt cổ chai lớn nhất, với độ trễ có thể lên tới hàng trăm mili-giây \cite{kawshan2026digital} [9].
    \item \textbf{Tính thiếu linh hoạt của Middleware:} Các giao thức truyền thông truyền thống (như TCP/IP hay UDP cơ bản trong ROS1) thường không được tối ưu để truyền đồng thời dữ liệu động học (kinematics) tần số cao và luồng video/Point Cloud chất lượng cao, gây ra hiện tượng nghẽn mạng \cite{luu2025enhancing} [10].
    \item \textbf{Sự lỗi thời của thiết bị tương tác:} Nhiều nghiên cứu trước đây bị giới hạn bởi phần cứng VR cũ (góc nhìn hẹp, thiếu khả năng Pass-through/Mixed Reality), làm giảm khả năng nhận thức môi trường thực tế (situational awareness) của người vận hành \cite{chen2023comparing} [11].
\end{itemize}

\section{Mục tiêu và Đóng góp của luận văn}
\label{sec:contributions_outline}
Luận văn này hướng tới việc nghiên cứu, thiết kế và phát triển một hệ thống điều khiển từ xa cho cánh tay robot công nghiệp (Denso VS-6577) sử dụng nền tảng ROS2, kết hợp đồng bộ thời gian thực với kính thực tế ảo Meta Quest 3 thông qua Unity 6. Các đóng góp khoa học và kỹ thuật (Contributions) của luận văn bao gồm:
\begin{enumerate}
    \item \textbf{Kiến trúc phần mềm phân tán dựa trên ROS2 và Ros2-Control:} Đề xuất và hiện thực một kiến trúc điều khiển tách biệt (decoupled), tận dụng ROS2 để giải quyết bài toán thời gian thực, khắc phục điểm nghẽn của ROS1 trong việc xử lý tín hiệu điều khiển tần số cao \cite{kawshan2026digital}.
    \item \textbf{Giao diện Digital Twin nhập vai trên Meta Quest 3:} Tích hợp môi trường Unity 6 và ROS-TCP-Connector để xây dựng Bản sao số đồng bộ theo thời gian thực. Hệ thống cung cấp phản hồi hình ảnh camera đa luồng và tương tác điều khiển trực quan (nắm/kéo/thả mô hình 3D), giảm thiểu tải nhận thức cho người dùng \cite{zhang2023research}.
    \item \textbf{Giải pháp tích hợp và tối ưu hóa điều khiển phần cứng đặc thù:} Hiện thực thành công chuỗi điều khiển kín (closed-loop) kết nối từ môi trường VR tới bộ điều khiển công nghiệp truyền thống (RC5). Luận văn đã xây dựng các module xử lý dữ liệu và bộ đệm trung gian nhằm khắc phục các giới hạn về độ trễ cơ học đặc thù của cánh tay robot, đảm bảo quỹ đạo do MoveIt2 lập kế hoạch được thực thi mượt mà và an toàn.
\end{enumerate}

% ---------------------------------------------------------

\chapter{CÁC NGHIÊN CỨU LIÊN QUAN}
\label{chap:related_works}
Trong những năm gần đây, sự hội tụ giữa Robot học (Robotics), Thực tế ảo (VR) và Bản sao số (Digital Twin) đã mở ra một kỷ nguyên mới cho kỹ thuật điều khiển từ xa. Chương này sẽ tiến hành phân tích sâu các tài liệu nghiên cứu hiện có, nhận diện những đặc trưng, cách làm chung và chỉ ra các khoảng trống nghiên cứu (Research Gaps) mà luận văn hiện tại sẽ cải thiện.

\section{Giao diện điều khiển và tương tác từ xa (Control Interfaces)}
Một trong những mục tiêu cốt lõi của teleoperation là làm cho việc điều khiển robot trở nên tự nhiên nhất với con người. Chen và cộng sự \cite{chen2023comparing} [4] đã thực hiện nghiên cứu so sánh ba phương thức: giao diện đồ họa 2D (GUI), nhận diện cử chỉ tay (hand gestures) và bộ điều khiển VR (VR controllers) trên cánh tay JAKA Minicobo. Kết quả chỉ ra rằng VR Controllers mang lại tốc độ hoàn thành công việc nhanh nhất, độ chính xác cao hơn và ít đòi hỏi sức lực tinh thần hơn so với GUI truyền thống \cite{chen2023comparing} [14, 15]. Tương tự, hệ thống BEAVR của MIT \cite{arclab2025beavr} [16] sử dụng Meta Quest 3S đã chứng minh rằng tay cầm VR cung cấp khả năng điều khiển linh hoạt cho các thao tác phức tạp (lên đến 90Hz) với độ trễ chỉ khoảng 10ms trong mạng nội bộ \cite{arclab2025beavr} [17, 18]. Tuy nhiên, các phương pháp này đòi hỏi sự đồng bộ rất cao về động học nghịch (Inverse Kinematics) để biến đổi tọa độ từ tay người sang không gian của robot.

\section{Tích hợp Bản sao số (Digital Twin) và Nhận thức không gian}
Để khắc phục nhược điểm của việc chỉ sử dụng camera 2D, công nghệ Digital Twin trong VR đã được ứng dụng mạnh mẽ. Hệ thống \textit{Vicarios} do Naceri và cộng sự \cite{naceri2021vicarios} [19] đề xuất đã tích hợp Unreal Engine và ROS để tạo ra một Bản sao số. Họ giới thiệu tính năng "teleporting" (dịch chuyển tức thời), cho phép người dùng thay đổi góc nhìn linh hoạt trong không gian 3D để tránh các góc khuất do cánh tay robot tạo ra \cite{naceri2021vicarios} [20, 21]. 

Ở mức độ cao hơn, Torrejón và cộng sự \cite{torrejon2026teleoperation} [22] đã kết hợp Digital Twin với dữ liệu Đám mây điểm (Point Cloud) mật độ cao (hơn 1 triệu điểm) để tạo ra một không gian làm việc bao gồm cả robot và môi trường từ xa. Dù mang lại sự chân thực cao, phương pháp này đòi hỏi năng lực xử lý đồ họa (GPU) cực lớn và gây ra độ trễ tổng thể xấp xỉ 190ms \cite{torrejon2026teleoperation} [6, 23]. Trong một hướng tiếp cận khác, Su và cộng sự \cite{su2022mixed} [24] sử dụng không gian Thực tế hỗn hợp (Mixed Reality) để lồng ghép không gian làm việc của người dùng và robot, giúp giảm thời gian hoàn thành nhiệm vụ xuống 48\% so với phương pháp 2D truyền thống \cite{su2022mixed} [25]. Những nghiên cứu này đều có một điểm chung: sử dụng Bản sao số như một vùng đệm (buffer) để người dùng thao tác không bị ảnh hưởng bởi độ trễ vật lý của mạng \cite{autolab2021telesurgery} [7, 26].

\section{Kiến trúc truyền thông và Middleware (Communication \& Middleware)}
Việc lựa chọn Middleware và giao thức mạng quyết định trực tiếp đến độ trễ (latency) và tính thời gian thực (real-time). 
\begin{itemize}
    \item \textbf{Phân tích nền tảng mô phỏng:} Singh và cộng sự \cite{singh2025comparative} [27] đã thực hiện một so sánh thực nghiệm giữa Unity và Gazebo trong việc xây dựng Bản sao số. Kết quả khẳng định Unity vượt trội về chất lượng đồ họa và độ trễ giao tiếp thấp (khoảng 77.67 ms) khi đồng bộ vị trí khớp, trong khi Gazebo mạnh hơn về khả năng tích hợp ROS gốc \cite{singh2025comparative} [28, 29].
    \item \textbf{Giao thức mạng:} Đối với các mạng công nghiệp (IIoT), Luu và cộng sự \cite{luu2025enhancing} [30] đề xuất sử dụng giao thức MQTT kết hợp với VR. Dù MQTT giúp duy trì kết nối ổn định khi rớt mạng, độ trễ trên diện rộng vẫn lên tới 99ms - 158ms tùy thuộc vào mức QoS \cite{luu2025enhancing} [31, 32]. 
    \item \textbf{Sự chuyển dịch sang ROS2:} Hầu hết các hệ thống trước đây đều dựa vào cơ chế Master/Slave của ROS1, vốn dễ gặp nút thắt cổ chai khi băng thông cao \cite{zhang2023research} [13]. Ali và cộng sự \cite{ali2025digital} [33] đã tiên phong áp dụng ROS2 kết hợp với thuật toán Soft Actor-Critic (SAC) trong in 3D, sử dụng ROS-TCP-Connector để liên kết Unity và phần cứng vật lý (Viper X300s). Nền tảng ROS2 dựa trên DDS đã giúp họ đạt được độ trễ đồng bộ gần như tức thời (khoảng 20ms) trong mạng cục bộ \cite{ali2025digital} [34]. 
\end{itemize}

\section{Phân tích khoảng trống nghiên cứu (Research Gaps)}
Thông qua việc phân tích kỹ lưỡng các tài liệu, luận văn nhận diện 2 khoảng trống nghiên cứu (Research Gaps) chính trong mảng điều khiển robot qua VR mà hệ thống hiện tại có thể đóng góp và cải thiện:
\begin{enumerate}
    \item \textbf{Gap 1: Tính thiếu linh hoạt và điểm nghẽn của ROS1 trong truyền thông đa luồng.} Đa số các ứng dụng HRC trực quan \cite{hetrick2020comparing, chen2023comparing} và các hệ thống Digital Twin hiện hành \cite{kawshan2026digital} vẫn sử dụng ROS1. Việc thiếu vắng một kiến trúc toàn diện dựa trên \textit{ROS2-Control} tích hợp \textit{MoveIt2} khiến các hệ thống cũ khó mở rộng, dễ gặp độ trễ lớn khi phải truyền tải đồng thời dữ liệu động học tần số cao và luồng hình ảnh/Point Cloud.
    \item \textbf{Gap 2: Giới hạn về nhận thức không gian do phần cứng VR thế hệ trước.} Nhiều nghiên cứu trước đây thường dùng Oculus Quest 2 hoặc HTC Vive \cite{naceri2021vicarios, chen2023comparing}, vốn có giới hạn về độ phân giải, độ trễ tracking và khả năng hiển thị xuyên thấu (Pass-through). Điều này làm giảm chất lượng của trải nghiệm Digital Twin. Luận văn tận dụng phần cứng hiện đại \textit{Meta Quest 3} kết hợp Unity 6 để mang lại trải nghiệm thực tế hỗn hợp (MR) sắc nét hơn, cải thiện nhận thức không gian mà không làm tăng độ trễ hệ thống.
\end{enumerate}

\section{Bảng thống kê các nghiên cứu liên quan}
Bảng \ref{tab:literature_review} trình bày thống kê tổng hợp các nghiên cứu nổi bật, cách tiếp cận của họ và điểm khác biệt của luận văn này nhằm lấp đầy các khoảng trống công nghệ đã nêu.

\begin{table}[H]
\centering
\caption{Bảng thống kê và so sánh các nghiên cứu liên quan về VR Teleoperation và Digital Twin}
\label{tab:literature_review}
\resizebox{\textwidth}{!}{%
\begin{tabular}{|p{3cm}|p{2.5cm}|p{2.5cm}|p{3.5cm}|p{3.5cm}|p{3.5cm}|}
\hline
\textbf{Tác giả / Năm} & \textbf{Robot Hardware} & \textbf{Middleware \& Truyền thông} & \textbf{Thiết bị VR/MR} & \textbf{Đặc trưng \& Phương pháp} & \textbf{Hạn chế (Gaps) / Khác biệt của luận văn} \\ \hline
Chen et al. (2023) \cite{chen2023comparing} & JAKA Minicobo & ROS, UDP & Oculus Quest 2 & So sánh VR Controller, Gesture, và GUI 2D. Dùng camera RGB. & Phụ thuộc ROS1, kính Oculus 2 cũ. Chưa giải quyết tốt độ nhiễu tín hiệu tay. \\ \hline
Naceri et al. (2021) \cite{naceri2021vicarios} & UR5, SoftHand & ROS, Unreal, ROSBridge & HTC Vive Pro & Giao diện "Vicarios", tính năng Teleporting đổi góc nhìn, stream Point Cloud. & ROS1, độ trễ Point Cloud cao, setup nhiều camera tĩnh phức tạp. \\ \hline
Luu et al. (2025) \cite{luu2025enhancing} & ABB IRB1200 & Giao thức MQTT (IIoT), SteamVR & HTC Vive & Đánh giá mạng IIoT với QoS khác nhau. Đo lường độ trễ mạng diện rộng (WAN). & Chỉ tập trung vào tối ưu mạng MQTT, chưa khai thác sâu tính thời gian thực của ROS2. \\ \hline
Torrejón et al. (2026) \cite{torrejon2026teleoperation} & Dual-arm (P3bot) & RoboComp/Ice, Unreal 5.6 & Meta Quest 3 & Render Point Cloud >1M điểm qua GPU (Niagara). Digital Twin đồng bộ cao. & Đòi hỏi cấu hình PC cực mạnh (RTX 3070+), độ trễ tổng thể sát ngưỡng (190ms). \\ \hline
Kawshan \& Peng (2026) \cite{kawshan2026digital} & JetCobot (7-DOF) & ROS1, ROS-TCP, Unity & Desktop/VR & Tối ưu hóa RRTConnect, phân tích độ trễ chi tiết (144ms end-to-end). & Chủ yếu đo lường trên ROS1, chưa áp dụng cơ chế Ros2-Control cho teleoperation. \\ \hline
\textbf{Luận văn hiện tại} & \textbf{Denso VS-6577} & \textbf{ROS2 (DDS), ROS-Sharp, UART} & \textbf{Meta Quest 3 (Unity 6)} & \textbf{Kiến trúc tách biệt (decoupled) với Ros2-Control \& MoveIt2; đồng bộ Digital Twin thời gian thực.} & \textbf{Khắc phục điểm nghẽn ROS1 bằng ROS2; tối ưu hóa tương tác trực quan với phần cứng VR mới nhất.} \\ \hline
\end{tabular}%
}
\end{table}

Đúc kết lại, dựa trên việc kế thừa các phương pháp xây dựng Bản sao số và phân tích độ trễ từ các nghiên cứu đi trước, luận văn sẽ khai thác triệt để sức mạnh của middleware \textbf{ROS2} kết hợp \textbf{Unity 6} và thiết bị \textbf{Meta Quest 3}. Sự khác biệt lớn nhất nằm ở việc chuyển dịch thành công từ kiến trúc ROS1 nguyên bản sang chuỗi điều khiển kín (closed-loop) ứng dụng tiêu chuẩn \textit{Ros2-Control}, tạo ra một nền tảng điều khiển từ xa có khả năng mở rộng cao và độ trễ thấp.

